\documentclass[a5paper, 10pt]{article}


\usepackage{cite, amsmath, amssymb, booktabs, lastpage, graphicx, float}
\usepackage[hidelinks]{hyperref}

\usepackage{geometry}
\geometry{margin=1.5cm,a5paper,twoside,inner=2cm}

\usepackage{setspace}


\usepackage{amsthm}
%\usepackage[
%			type={CC},
%			modifier={by},
%			version={4.0},
%			imagewidth=5em,
%			imagedistance=5em,
%			]{doclicense}
%
%\makeatletter
%	\def\blfootnote{\gdef\@thefnmark{}\@footnotetext}
%\makeatother

\theoremstyle{plain}
\newtheorem{lemma}{Lemma}
\newtheorem{theorem}{Theorem}
\newtheorem{conjecture}{Conjecture}
\newtheorem{proposition}[theorem]{Proposition}
\newtheorem*{corollary}{Corollary}
\renewcommand{\qedsymbol}{\rule{.7em}{.7em}}

\theoremstyle{remark}
\newtheorem*{remark}{Remark}
\newtheorem*{note}{Note}
\newtheorem{case}{Case}

\theoremstyle{definition}
\newtheorem{definition}{Definition}
\newtheorem{example}{Example}

%% Changing the title, author, etc.
\usepackage{graphicx}
\makeatletter
\renewcommand{\maketitle}{
    \begin{center}
    \rule[.2em]{\textwidth}{0.353mm}
        \begin{minipage}[m]{0.35\textwidth}
            {\scriptsize
                \begin{center}
                    \textsf{\textbf{\huge MATCH}\\
                        \textit{Communications in Mathematical\\
                            and in Computer Chemistry}
                    }
            \end{center}}
        \end{minipage}\hfill
        \begin{minipage}[m]{0.65\textwidth}
            \begin{flushright}
                \baselineskip=10pt
                {\scriptsize\sffamily{\itshape  MATCH Commun. Math. Comput. Chem.}
                \textbf{\vol{}}
                (\pubyear)
                    \thepage--\pageref{LastPage}}\\
                {\scriptsize\sffamily {\bfseries ISSN:} 0340--6253}\\
                {\scriptsize\sffamily \textbf{doi:} \doi{10.46793/match.\vol-\num.\id}}
            \end{flushright}
        \end{minipage}
        \rule[1em]{\textwidth}{.353mm}
        \baselineskip=0.30in
        {\Large\bfseries \@title} \par
        \vspace{5mm}
        \baselineskip=0.2in
        {\large\bfseries \@author}\par
        \vspace{1mm}
        {\it \@address} \par
        {\small\tt \@email} \par
        \vspace{3mm}
        {\small (Received \@date)} \par
    \end{center}
    \vspace{3mm}
}
\newcommand{\address}[1]{\def\@address{#1}}
\newcommand{\email}[1]{\def\@email{#1}}
\newcommand{\doi}[1]{\href{https://doi.org/#1}{#1}}

\newcommand{\vol}{00}
\newcommand{\num}{00}
\newcommand{\pubyear}{0000}
\newcommand{\id}{xxxxx}
%%setting page-number
\setcounter{page}{1}


\thispagestyle{empty}
\makeatother

%% Making << \acknowledgment >> command
\newcommand{\acknowledgment}[1]{\vspace{5mm}\singlespacing
    {\noindent\textbf{\textit{Acknowledgment\/}:} #1}
}

%% Setting page numbers
\usepackage{fancyhdr}

%% now reset the headers and footers
\fancyhead{}
\fancyfoot{}
\renewcommand{\headrulewidth}{.3pt}
\voffset = 18pt
\headsep = 3pt

\renewcommand*{\thefootnote}{\fnsymbol{footnote}}

\usepackage[margin=1cm,%
            font=footnotesize,%
            format=hang,%
            labelsep=period,%
            labelfont=bf]{caption}

%%%%%%%%%%%%%%%%%%%%%%%%%%%%%%%%%%%%%%%%%%%%%%%%%%
% PLACE ADDITIONAL PACKAGES HERE (if it is needed):
\usepackage[T1]{fontenc}
\usepackage{lmodern}
\usepackage{mathtools}
\usepackage{array}
\usepackage{enumitem}
\setlist[itemize]{topsep=2pt,itemsep=2pt}
\setlist[enumerate]{topsep=2pt,itemsep=2pt}
\emergencystretch=3em
\hbadness=2000
\setlength{\marginparwidth}{2cm}
%%%%%%%%%%%%%%%%%%%%%%%%%%%%%%%%%%%%%%%%%%%%%%%%%%

%% make the odd pages have the section name on the top right
\fancyhead[RO]{\footnotesize \thepage}
%% make the even pages have the chapter name on the top left
\fancyhead[LE]{\footnotesize \thepage}

%% making page-numbers visible
\pagestyle{fancy}

% Custom macros for the Loyola index family
    \newcommand{\Wuzi}{LO(G;\alpha,\beta,\gamma)}
  \newcommand{\edgepsi}{\psi(x,y;\alpha,\beta,\gamma)}
  \newcommand{\bid}{\mathcal{F}_{\mathrm{BID}}}
% Path to figures
\graphicspath{{figures/}{./}{../figures/}}

% robust figure inclusion: render if present (root or figures/), else a small note
\newcommand{\safefig}[2]{%
  \IfFileExists{#1}{\includegraphics[width=#2\linewidth]{#1}}{%
   \IfFileExists{figures/#1}{\includegraphics[width=#2\linewidth]{#1}}{%
    \fbox{\parbox[c][3cm][c]{0.8\linewidth}{\centering Figure not in project --- upload \texttt{\detokenize{#1}}.}}}}%
}

%%%%%%%%%%%%%%%%%%%%%%%%%%%%%%%%%%%%%%%%%%%%%%%%%%%%%%%%%%%%%%%%%%%%%%%%%%%%%%%%%%%%%%%

\title{The Loyola Index Family: What a Degree-Imbalance Coordinate Can and Cannot Resolve}
\author{Advik Natarajan$^{a}$, Ganesh Shiwakoti$^{b}$, Michael Arockiaraj$^{a,}$\footnote{Corresponding author.}}

\address{$^a$Department of Mathematics, Loyola College, Chennai, India\\
    $^b$Department of Mathematics and Statistics, Florida Atlantic University, USA
    }

\email{marockiaraj@gmail.com}

\date{\today}

% !TeX program = pdflatex

\begin{document}

\maketitle

\begin{abstract}%\blfootnote{\doclicenseThis}
  The two-parameter Gutman--Milovanovi\'c family
  $M_{\alpha,\beta}(G) = \sum_{uv \in E(G)} (d_u d_v)^{\alpha} (d_u + d_v)^{\beta}$
  unifies the Zagreb, Randi\'c, sum-connectivity and related degree-based
  indices. We study its normalized degree-imbalance exponential
  extension, the \emph{Loyola index family}
  \begin{equation*}
  LO(G;\alpha,\beta,\gamma)
  \;=\; \sum_{uv \in E(G)}
  (d_u d_v)^{\alpha}\,(d_u + d_v)^{\beta}\,
  \exp\!\left(\gamma\,\frac{|d_u - d_v|}{d_u + d_v}\right),
  \end{equation*}
  with $LO(G;\alpha,\beta,0) = M_{\alpha,\beta}(G)$ identically.
  We ask what the imbalance coordinate $\gamma$ adds to the parent
  family. On a ratio plane that separates size from shape, the
  pure-$\gamma$ points distinguish two graphs exactly when their
  histograms of edge imbalances differ. On small trees this is the
  discrimination limit of all degree-based indices, and we prove that
  the limit is lost at larger orders. We bound the number of times two
  graphs exchange order along $\gamma$ and characterise the directions
  in parameter space that a given graph can identify. Pre-registered
  tests on alkane properties show that $\gamma$ gives no predictive gain
  over the parent family.
  \end{abstract}

\onehalfspacing

\section{Introduction}\label{sec:intro}
  % =====================================================================

  Degree-based topological indices form one of the most developed
  descriptor classes of chemical graph theory. The class runs from the
  Randi\'c and Zagreb indices~\cite{Randic1975, GutmanTrinajstic1972}
  and their parametric extensions, the general Randi\'c and
  sum-connectivity families~\cite{BollobasErdos1998, ZhouTrinajstic2010},
  to the Sombor family and its many variants. See the
  survey~\cite{Gutman2026SomborSurvey}. Extremal results on
  bond-incident-degree (BID) indices are surveyed
  in~\cite{AliGutmanFurtula2026BID}, while a finite-dimensional linear
  framework for such indices is developed in~\cite{Rada2026LinearVDB}.

  The \emph{Gutman--Milovanovi\'c family}, introduced
  in~\cite{GutmanMilovanovic2020AKCE} and studied in the recent bounds
  and QSPR literature
  \cite{GranadosJMC2025GM, MolinaSigarreta2025GM, PortillaJMC2025GM},
  unifies the product-based and sum-based degree indices with two
  parameters. It is defined by $M_{\alpha,\beta}(G) = \sum_{uv \in E(G)}
  (d_u d_v)^{\alpha} (d_u + d_v)^{\beta}$, $\alpha, \beta \in
  \mathbb{R}$. Its special cases are catalogued
  in~\cite{MolinaSigarreta2025GM}. They include the first and second
  Zagreb indices $M_{0,1} = M_1$ and $M_{1,0} = M_2$, the
  hyper-Zagreb index $M_{0,2} = HM$~\cite{ShirdelRezapourSayadi2013HM},
  the modified second Zagreb index
  $M_{-1,0} = {}^{m}\!M_2$ and the Randi\'c index $M_{-1/2,0} = R$.
  Further special cases are the sum-connectivity index
  $M_{0,-1/2} = \chi$~\cite{ZhouTrinajstic2009sumconn}, the harmonic
  index $2M_{0,-1} = H$, the inverse sum indeg index $M_{1,-1} = ISI$,
  the geometric--arithmetic index
  $2M_{1/2,-1} = GA$~\cite{VukicevicFurtula2009GA} and the
  arithmetic--geometric index $\tfrac12 M_{-1/2,1} = AG$. The general
  Randi\'c family $R_\alpha = M_{\alpha,0}$ and the general
  sum-connectivity family $\chi_\beta = M_{0,\beta}$ are its two
  parameter axes, and parameter-optimised versions of the family have
  already been fitted to octane
  properties~\cite{AliGutmanFurtulaAlbalahiHamza2025}. The family,
  its special cases and its tunability are therefore not new.

  In this paper we study a three-parameter extension of
  $M_{\alpha,\beta}$, the \emph{Loyola index family}, obtained by
  attaching a separately parameterised exponential factor in the
  normalised degree imbalance $q_{uv} = |d_u - d_v|/(d_u + d_v)$ of
  each edge. The definition is $LO(G;\alpha,\beta,\gamma) = \sum_{uv}
  (d_u d_v)^{\alpha}(d_u+d_v)^{\beta} e^{\gamma q_{uv}}$, so that
  $LO(G;\alpha,\beta,0) = M_{\alpha,\beta}(G)$ identically, and positive
  $\gamma$ exponentially upweights degree-imbalanced edges while
  negative $\gamma$ downweights them. On regular graphs $\gamma$ has no
  effect. The paper therefore concerns $\gamma$ alone. We ask whether
  the imbalance direction is independent of the parent parameters, what
  analytic structure it carries, and whether it adds structural or
  predictive information beyond an optimised Gutman--Milovanovi\'c
  parent.

  Writing the summand as
  $\exp\bigl(\alpha \ln(d_u d_v) + \beta \ln(d_u + d_v) +
  \gamma\, q_{uv}\bigr)$ shows that the family is a three-parameter
  subfamily of the exponential vertex-degree-based indices of
  Rada~\cite{Rada2019ExpVDB}, whose extremal theory is developed
  in~\cite{Sigarreta2022ExpVDB, GaoAMC2024, HuMATCH2022}. Extremal
  problems are not considered here. Published one-parameter
  exponential indices (e.g.~\cite{DasMondal2024ExpGA,
  DasMondalRazaPal2025ExpISDD}) place their parameter on a single
  classical generator inside the exponent. Here the two
  Gutman--Milovanovi\'c basis terms act as the factors
  $(d_u d_v)^{\alpha}(d_u + d_v)^{\beta}$, while the imbalance
  coordinate carries its own parameter $\gamma$. The imbalance ratio
  itself is not new. It appears in the diminished Sombor
  index~\cite{MovahediGutmanRedzepovicFurtula2026} and in the
  irregularity index
  $IRLA(G) = 2\sum_{uv} q_{uv}$~\cite{RetiSharafdini2018IRLA}. The
  pure-$\gamma$ point $LO(G;0,0,\gamma) = \sum_{uv} e^{\gamma
  q_{uv}}$ is the exponential index of the latter in the sense
  of~\cite{Rada2019ExpVDB}. The ratio also appears in degree-ratio
  indices~\cite{NyauliBuragohain2027DRSO, GutmanBuragohainNyauli2026}
  and in the parent family itself, since
  $1 - q_{uv}^2 = 4 d_u d_v/(d_u + d_v)^2$ gives
  $M_{\alpha,-2\alpha}(G) = 4^{-\alpha} \sum_{uv} (1 - q_{uv}^2)^{\alpha}$.
  What is new is a coordinate \emph{linear} in $q_{uv}$ in the
  exponent. Unlike $\ln(1 - q_{uv}^2)$, this coordinate is not an
  affine function of $\ln(d_u d_v)$ and $\ln(d_u + d_v)$ once a vertex
  of degree at least $3$ is admitted. We prove this below. To our
  knowledge, and within the scope of the literature search documented
  in the reproducibility package, no published index couples it
  multiplicatively, with its own parameter, to the
  Gutman--Milovanovi\'c family. The exponential form is not derived
  from a physical model and is chosen instead because it makes the
  family log-linear in its parameters. Then $\gamma$ acts
  multiplicatively and only on irregularity, and the whole family is
  an exponential family in $(\alpha,\beta,\gamma)$, which makes
  available the convexity, identifiability and discrimination results
  that follow.

  On hydrogen-suppressed molecular graphs with $\Delta \le 4$, every
  BID index is a linear functional of the ten edge-degree-pair counts
  $m_{ij}(G) = |\{uv \in E(G) : \{d_u, d_v\} = \{i, j\}\}|$,
  $1 \le i \le j \le 4$ (compare the linear framework of
  Rada~\cite{Rada2026LinearVDB}). Hence the Loyola family lies inside a
  fixed ten-dimensional descriptor space, however its parameters are
  chosen.

  Several of the mathematical results apply standard tools, namely
  log-sum-exp convexity, Laguerre's rule of signs, the
  Lindemann--Weierstrass theorem and exponential-family minimality.
  The contribution lies in the structural characterisation they yield.
  The family is not proposed as a better predictor, and we find no
  predictive advantage on the sets and protocols studied here. Its
  interest lies in the structure it exposes. It keeps the interpretation
  of the parent family, separates size from shape, and shares the
  ratio-plane geometry of $GA$ and $AG$. Because the QSPR datasets are
  small ($18$ octanes, $35$ nonanes, $34$ decanes, $100$ pooled C6--C10
  alkanes), the chemometric results are diagnostic evidence rather
  than proof of chemical predictive utility.

  The contributions are as follows. (i) The scale--shape factorisation
  and the ratio-plane theorem (Theorem~\ref{thm:ratioplane}). This
  theorem characterises exactly what the pure-$\gamma$ line can
  discriminate and is in large part a limitation result, with tree
  collisions at every order from $13$, and from $16$ for
  $\Delta \le 4$. (ii) The crossing bound
  (Proposition~\ref{prop:laguerre}, from Laguerre's rule of signs) and
  per-graph identifiability (Proposition~\ref{prop:identrank}), the
  latter resting on an elementary independence lemma
  (Lemma~\ref{prop:indep}). (iii) Log-convexity with mean/covariance
  derivatives (Theorem~\ref{thm:logconvex}, standard), and closed forms
  and routine termwise bounds for $\gamma \ge 0$ (Supplementary
  Material). (iv) The degree-pair ceiling, an elementary observation
  that is realised on the octane benchmark
  (Proposition~\ref{prop:collision}), and structure sensitivity. (v) A
  negative, reproducible QSPR diagnostic, in which freeing $\gamma$
  gives no consistent gain once the parent's search box does not
  truncate it. This is complemented by pre-registered tests on measured
  boiling points of decanes and of pooled C6--C10 alkanes
  (Section~\ref{sec:bp}). There, ridge regression on the degree-pair
  counts is far ahead of every single index on the decanes and ahead on
  the pooled set.


% =====================================================================
\section{Preliminaries}\label{sec:prelim}
% =====================================================================

\subsection{Notation}\label{sec:notation}

Let $G = (V, E)$ be a finite simple connected graph of order
$n = |V|$ and size $m = |E|$, and for a vertex $v \in V$ let the degree
$d_v = d_G(v)$ count the edges incident to $v$. The maximum and
minimum degrees are denoted $\Delta = \Delta(G)$ and
$\delta = \delta(G)$. The chemometric sections use hydrogen-suppressed
molecular graphs of alkane carbon skeletons, for which $\Delta \le 4$,
while the mathematical results do not assume this bound unless stated.

For an edge $e = uv \in E(G)$ we write $(d_u, d_v)$ for its
\emph{degree pair}, understood as an unordered multiset
$\{d_u, d_v\}$.


\subsection{Definition}\label{sec:def}

\begin{definition}[Loyola index family]\label{def:wuzi}
Let $G = (V, E)$ be a simple connected graph. For
$\alpha, \beta, \gamma \in \mathbb{R}$, the \emph{Loyola index family}
is defined as,
\begin{equation}\label{eq:wuzi}
\Wuzi
\;=\; \sum_{uv \in E(G)}
(d_u d_v)^{\alpha}\,(d_u + d_v)^{\beta}\,
\exp\!\left(\gamma\,\frac{|d_u - d_v|}{d_u + d_v}\right)
\end{equation}
\end{definition}

Eq.~\eqref{eq:wuzi} is well defined for all
$\alpha, \beta \in \mathbb{R}$ on graphs without isolated vertices.
Each summand is symmetric in $(u,v)$, so $\Wuzi$ depends only on the
multiset of unordered edge degree pairs of~$G$. When $\gamma = 0$
the family reduces identically to the Gutman--Milovanovi\'c
family~\cite{GutmanMilovanovic2020AKCE, MolinaSigarreta2025GM},
\begin{equation*}
LO(G;\alpha,\beta,0) \;=\; M_{\alpha,\beta}(G)
\;=\; \sum_{uv \in E(G)} (d_u d_v)^{\alpha} (d_u + d_v)^{\beta},
\end{equation*}
whose classical special cases are catalogued in the Introduction.
Throughout, $q_{uv} = \mathrm{imb}(d_u, d_v)$, where
$\mathrm{imb}(x, y) = |x - y|/(x + y)$, denotes the
normalised degree imbalance of the edge $uv$, and positive $\gamma$
exponentially upweights degree-imbalanced edges while negative $\gamma$
downweights them. On regular graphs every $q_{uv} = 0$, so $\gamma$
has no effect and
$LO(G;\alpha,\beta,\gamma) = M_{\alpha,\beta}(G)$ for all $\gamma$.
 % =====================================================================
  \section{Parameter geometry}\label{sec:geometry}

  Throughout,
  $z_{uv} = \bigl(\ln(d_u d_v),\, \ln(d_u + d_v),\, q_{uv}\bigr)$ is
  the feature vector of the edge $uv$ and
  $\theta = (\alpha, \beta, \gamma)$, so that
  $LO_G(\theta) = \sum_{uv \in E(G)} e^{\theta^{\top} z_{uv}}$.


  \begin{lemma}[Affine independence of the imbalance
  direction]\label{prop:indep}
  For $\Delta \ge 3$, with
  $P_{\Delta} = \{(i,j) : 1 \le i \le j \le \Delta\}$ the
  admissible degree pairs, the augmented feature vectors
  $\bigl(1, \ln(ij), \ln(i+j), q_{ij}\bigr)$, $(i,j) \in P_{\Delta}$,
  span $\mathbb{R}^4$. Consequently, if two parameter triples
  $\theta, \theta'$ satisfy
  $e^{\theta^{\top} z} = c\, e^{\theta'^{\top} z}$ for some constant
  $c > 0$ and all $(i,j) \in P_{\Delta}$, then $\theta = \theta'$ and
  $c = 1$. In particular, no choice of $(\alpha', \beta')$ and scaling
  constant reproduces the action of $\gamma \neq 0$ on all admissible
  degree pairs. For $\Delta = 2$ the span has dimension $3$, while for
  $\Delta = 1$ the only pair is $(1,1)$, with augmented vector
  $(1, 0, \ln 2, 0)$, and the span has dimension $1$. In both cases the
  conclusion fails.
  \end{lemma}

  \begin{proof}
  It suffices to show that $q_{ij}$ is not an affine function of
  $\bigl(\ln(ij), \ln(i+j)\bigr)$ on $P_{\Delta}$, and the four pairs
  $(1,1), (2,2), (1,2), (1,3)$, all admissible once $\Delta \ge 3$,
  witness this. Suppose
  $q_{ij} = a + b \ln(ij) + c \ln(i+j)$ on these pairs. The pairs
  $(1,1)$ and $(2,2)$ have $q = 0$, so they give $a + c \ln 2 = 0$ and
  $a + (2b + 2c) \ln 2 = 0$, hence $a = -c \ln 2$ and $b = -c/2$.
  Substituting into the $(1,2)$ and $(1,3)$ equations gives
  $\tfrac13 = c\,(\ln 3 - \tfrac32 \ln 2)$ and
  $\tfrac12 = c\,(\ln 2 - \tfrac12 \ln 3)$. In particular $c \neq 0$.
  Three times the first minus twice the second gives
  $0 = c\,(4 \ln 3 - \tfrac{13}{2} \ln 2)$, hence $8 \ln 3 = 13 \ln 2$,
  i.e.\ $3^{8} = 2^{13}$. This is false ($6561 \neq 8192$). Hence the
  augmented vectors have rank $4$. If
  $e^{\theta^{\top} z} = c\, e^{\theta'^{\top} z}$ on $P_{\Delta}$,
  then $(\theta - \theta')^{\top} z_{ij} = \ln c$ for all admissible
  pairs, so $(\ln c, \theta' - \theta)$ is orthogonal to every
  augmented vector, and $\theta = \theta'$ and $c = 1$. For
  $\Delta = 2$ only the pairs $(1,1), (1,2), (2,2)$ exist, and a
  direct computation shows the augmented vectors have rank $3$, while for
  $\Delta = 1$ the single augmented vector has rank $1$.
  \end{proof}

  \begin{remark}[Independence versus per-graph
  identifiability]
  Lemma~\ref{prop:indep} concerns the edge-weight
  \emph{function} on the full admissible pair set rather than
  every individual graph. On a single graph $G$ some direction $v$ satisfies
  $LO(G;\theta + tv) = e^{ct}\,LO(G;\theta)$ for all $t$, unless the
  realised feature vectors $z_{ij}$ affinely span $\mathbb{R}^3$
  (Proposition~\ref{prop:identrank}). Along such a direction the index
  only rescales. On a regular graph nothing about $\theta$
  is identifiable beyond one linear functional. If the edges of $G$
  realise at most three affinely independent augmented vectors, some
  direction of $\theta$, possibly the $\gamma$ direction, is unidentifiable
  from $LO(G;\cdot)$ alone. The lemma guarantees
  that no such degeneracy is forced on the whole molecular class
  $\Delta \ge 3$, so the family separates parameters as a family of
  descriptors, although a particular graph may not.
  \end{remark}


  The next result is the standard convexity of a log-sum-exp
  (log-partition) function, stated for $LO$ because strictness depends
  on the degree pairs that $G$ realises.

  \begin{theorem}[Global log-convexity]\label{thm:logconvex}
  Fix a graph $G$ with $m \ge 1$ edges, define
  $F_G(\theta) = \log LO_G(\theta)$ on $\mathbb{R}^3$, and let
  $\mathbb{E}_{\theta}$ and $\operatorname{Cov}_{\theta}$ denote mean
  and covariance under the edge probability measure
  $w_{uv}(\theta) = e^{\theta^{\top} z_{uv}} / LO_G(\theta)$. Then
  $F_G$ is convex, with
  \begin{equation*}
  \nabla F_G(\theta) = \mathbb{E}_{\theta}[z], \qquad
  \nabla^2 F_G(\theta) = \operatorname{Cov}_{\theta}(z) \succeq 0 .
  \end{equation*}
  In particular
  $\partial_{\gamma} F_G = \mathbb{E}_{\theta}[q] \in [0, 1)$ and
  $\partial^2_{\gamma} F_G = \operatorname{Var}_{\theta}(q) \ge 0$, so
  $\log LO$ is nondecreasing and convex in $\gamma$, and it is strictly
  increasing in $\gamma$ unless every edge of $G$ joins vertices of
  equal degree. Moreover
  $\operatorname{Var}_{\theta}(q) = 0$ at some (equivalently, every)
  $\theta$ if and only if $q_{uv}$ is constant over the edges of $G$,
  which covers regular graphs ($q \equiv 0$) and, for example,
  $(\Delta,\delta)$-biregular graphs. On these every edge carries the
  same nonzero imbalance, and $\log LO$ is affine and strictly
  increasing in $\gamma$. $F_G$ is strictly convex if and only if the
  edge feature vectors $\{z_{uv}\}_{uv \in E}$ do not lie in a proper
  affine subspace of $\mathbb{R}^3$. On a regular graph
  $F_G(\theta) = \log m + \theta^{\top} z_0$ is affine and
  $\operatorname{Cov}_{\theta}(z) = 0$.
  \end{theorem}

  The formulas follow by direct differentiation (proof in Supplementary
  Section~S9), and the Cauchy--Schwarz midpoint inequality there is midpoint
  log-convexity, a corollary of this theorem.

  \subsection{The ratio plane, crossings and identifiability}\label{sec:ratioplane}

  Since $4ij/(i+j)^2 = 1 - q_{ij}^2$, every admissible pair satisfies
  $\ln(ij) = 2\ln(i+j) + \ln(1 - q_{ij}^2) - \ln 4$, and therefore
  \begin{equation}\label{eq:factor}
  LO(G;\alpha,\beta,\gamma) = 4^{-\alpha} \sum_{uv \in E(G)}
  (d_u + d_v)^{2\alpha+\beta}\,(1 - q_{uv}^2)^{\alpha}\, e^{\gamma q_{uv}}.
  \end{equation}
  Thus $\kappa = 2\alpha + \beta$ acts only on the edge degree sum
  (scale), while the parameters $\alpha$, $\gamma$ act only on $q_{uv}$, a function of the
  degree ratio (shape). Write
  $n_q(G)$ for the number of edges of $G$ with $q_{uv} = q$ (the
  $q$-histogram) and $\Pi = \{(\alpha,\beta,\gamma) : \beta = -2\alpha\}$
  for the plane $\kappa = 0$, which contains every pure-$\gamma$ point
  $(0,0,\gamma)$, $GA = 2\,LO(\tfrac12,-1,0)$ and
  $AG = \tfrac12 LO(-\tfrac12,1,0)$.

  \begin{theorem}[Ratio plane]\label{thm:ratioplane}
  Let $G$ and $H$ be graphs.
  \begin{enumerate}
  \item[(i)] On $\Pi$,
  $LO(G;\alpha,-2\alpha,\gamma) = \sum_q n_q(G)\,\bigl((1-q^2)/4\bigr)^{\alpha}
  e^{\gamma q}$ depends on $G$ only through its $q$-histogram, and $G$ and
  $H$ agree at every point of $\Pi$ if and only if $n_q(G) = n_q(H)$
  for all $q$. Otherwise their values differ at every point of $\Pi$
  outside a set of measure zero.
  \item[(ii)] For every nonzero algebraic $\gamma$,
  $LO(G;0,0,\gamma) = LO(H;0,0,\gamma)$ if and only if
  $n_q(G) = n_q(H)$ for all $q$. On any finite set of graphs the same
  holds for all but finitely many real $\gamma$.
  \item[(iii)] Let $G$ and $H$ be trees with distinct
  edge-degree-pair profiles and equal $q$-histograms, each having a
  leaf adjacent to a vertex of the same degree $d \ge 2$. Attaching a
  pendant path of length $L \ge 1$ at such a leaf in each produces
  trees with the same two properties and the same maximum degree.
  Consequently, for every $n \ge 13$ (and every $n \ge 16$ among trees
  with $\Delta \le 4$) there are non-isomorphic trees of order $n$
  with distinct degree-pair profiles that no point of $\Pi$ separates.
  \end{enumerate}
  \end{theorem}

  \begin{proof}
  (i) is immediate from~\eqref{eq:factor}, since setting $\beta = -2\alpha$
  removes the degree-sum factor. For (ii) and (iii), agreement for equal
  histograms is immediate. Suppose the
  histograms differ. At $\alpha = 0$ the difference is
  $\sum_q (n_q(G) - n_q(H))\,e^{\gamma q}$, an exponential sum with
  distinct exponents and not all coefficients zero, so it is not
  identically zero. The difference on $\Pi$ is therefore a nonzero
  real-analytic function of $(\alpha,\gamma)$, whose zero set has
  measure zero. (ii) For algebraic $\gamma \neq 0$ the exponents
  $\gamma q$, $q \in \mathbb{Q}$, are distinct algebraic numbers, so by
  the Lindemann--Weierstrass theorem~\cite{Baker1975} the numbers
  $e^{\gamma q}$ are linearly independent over the algebraic numbers.
  Hence the integer coefficients $n_q(G) - n_q(H)$ must all vanish. For real
  $\gamma$, a nonzero exponential sum with $k$ distinct real exponents
  has at most $k - 1$ real zeros (Proposition~\ref{prop:laguerre}), and
  a finite set of graphs gives finitely many such sums. (iii) The
  attachment changes the profile of each tree by the same vector,
  $-(1,d) + (2,d) + (L-1)(2,2) + (1,2)$, so distinct profiles stay
  distinct and equal histograms stay equal. The new or changed
  vertices have degree at most $2$, so the maximum degree is unchanged
  when it is at least $2$. The order-$13$ pair of
  Section~\ref{sec:degeneracy} has a common leaf-neighbour degree
  $d = 3$, and the order-$16$ pair with $\Delta \le 4$,
  $\{(1,2)^{2}, (1,4)^{8}, (2,4)^{2}, (4,4)^{3}\}$ and
  $\{(1,4)^{8}, (2,2)^{3}, (2,4)^{4}\}$, has $d = 4$, which gives every
  larger order.
  \end{proof}

  Theorem~\ref{thm:ratioplane} is both a limitation result and a
  characterisation. By exhaustive enumeration of orders $7$--$17$ (orders
  below $7$ are trivial), $q$-histograms coincide with degree-pair
  profiles for all trees of order at most $12$. The same holds for trees with
  $\Delta \le 4$ of order at most $15$. So the pure-$\gamma$ points
  attain the degree-pair floor exactly up to those orders and, by part~(iii),
  at no larger order. The tie of $GA$, $AG$ and the
  pure-$\gamma$ points on the trees of order $10$
  (Section~\ref{sec:degeneracy}) is consistent with this, since all lie on
  $\Pi$. At orders $13$--$17$ generic points off
  $\Pi$ (Proposition~\ref{prop:generic}) and the Randi\'c index
  separate more trees than any point of $\Pi$. In the other direction,
  $GA$, whose values are sums of square roots, merges two distinct
  $q$-histograms among trees first at order $17$. A witness is the pair
  $\{(1,5)^{10}, (2,2)^{2}, (5,5)^{2}, (2,5), (1,2)\}$ and
  $\{(4,5)^{3}, (1,4)^{5}, (1,5)^{6}, (2,5), (1,2)\}$, which have equal $GA$
  but different $q$-histograms, while at that order the pure-$\gamma$ points,
  by part~(ii), still separate all histograms. That $GA$ depends only
  on the degree ratio is well known, and this paper adds the exact
  ceiling and the colliding families.

  \begin{proposition}[$\gamma$-crossings]\label{prop:laguerre}
  Fix $(\alpha,\beta)$ and graphs $G$, $H$, let
  $q_1 < \dots < q_r$ be the distinct imbalances of the degree pairs
  occurring in $G$ or $H$, and write
  $D(\gamma) = LO(G;\alpha,\beta,\gamma) - LO(H;\alpha,\beta,\gamma)
  = \sum_{k=1}^{r} c_k e^{\gamma q_k}$ with
  $c_k = \sum_{q_{ij} = q_k} (m_{ij}(G) - m_{ij}(H))(ij)^{\alpha}(i+j)^{\beta}$.
  Let $V$ be the number of sign changes in $(c_1,\dots,c_r)$ after
  zero terms are removed. If $D \not\equiv 0$, then $D$ has at most
  $V \le r - 1$ real zeros counted with multiplicity, and their number
  has the parity of $V$. In particular, for graphs with $\Delta \le 4$
  the possible imbalances are $0, \tfrac17, \tfrac15, \tfrac13,
  \tfrac12, \tfrac35$, and two such graphs exchange order at most five times as
  $\gamma$ ranges over $\mathbb{R}$, for every $(\alpha,\beta)$.
  \end{proposition}

  \begin{proof}
  With $x = e^{\gamma} > 0$, $D$ becomes the generalised polynomial
  $\sum_k c_k x^{q_k}$ with distinct real exponents, so Laguerre's
  extension of Descartes' rule of signs applies~\cite{Jameson2006}. The
  number of positive zeros, with multiplicity, is at most $V$ and has
  its parity, and the parity statement also follows from the signs of $D$
  as $\gamma \to \mp\infty$. These are the signs of the first and last nonzero $c_k$,
  by the termwise argument of Proposition~\ref{prop:gamma_eq_0_or_inf}
  applied to $D$.
  \end{proof}

  The bound $V$ can be attained. Two trees of order $12$ with
  maximum degree $6$ and eight distinct imbalances have $V = 5$ and cross
  five times at $(\alpha,\beta) = (-1,0)$, near $\gamma = -17.2$,
  $-12.0$, $-2.0$, $2.7$ and $31.1$. By the bound these are all its
  zeros. Whether five crossings are attained for $\Delta \le 4$ is open. Along directions other than the $\gamma$ axis the exponents are
not the $q_k$, so the constant five does not apply.

  The per-graph identifiability discussed in the remark after Lemma~\ref{prop:indep} can be made
  exact. Let $R(G)$ be the set of
  degree pairs realised by $G$, $r(G)$ the number of distinct imbalances
  among them, and $A(G)$ the affine dimension of
  $\{z_{ij} : (i,j) \in R(G)\}$.

  \begin{proposition}[Per-graph scale-invariant directions]\label{prop:identrank}
  \begin{enumerate}
  \item[(i)] $A(G)$ equals the rank of the Hessian
  $\operatorname{Cov}_{\theta}(z)$ of $\log LO_G$ at every $\theta$, and the
  directions $v$ with $LO(G;\theta + tv) = e^{ct}\,LO(G;\theta)$ for
  all $t$ and $\theta$ form a subspace of dimension $3 - A(G)$.
  \item[(ii)] $A(G) \le \min\bigl(|R(G)| - 1,\ r(G),\ 3\bigr)$ for every
  graph.
  \item[(iii)] If $r(G) = 2$ with imbalances $q_1 \neq q_2$, then
  $v = (1, -2, c)$ with
  $c = [\ln(1 - q_2^2) - \ln(1 - q_1^2)]/(q_1 - q_2)$ is such a
  direction.
  \end{enumerate}
  \end{proposition}

  \begin{proof}
  (i) is the standard minimal-representation property of exponential
  families~\cite{Brown1986}, and since $\operatorname{Cov}_{\theta}(z)$ has full
  support on $R(G)$, its rank is the affine dimension of the support.
  Also $v^{\top} z$ is constant on $R(G)$ exactly for $v$ orthogonal to
  the differences of the $z_{ij}$. (ii) By~\eqref{eq:factor},
  $z_{ij} = (2s + \ln(1-q^2) - \ln 4,\ s,\ q)$ with $s = \ln(i+j)$, so
  pairs of equal imbalance lie on one line parallel to $(2,1,0)$, and
  $r$ parallel lines span an affine space of dimension at most $r$.
  (iii) $v^{\top} z_{ij} = \ln(1 - q_{ij}^2) - \ln 4 + c\,q_{ij}$ takes
  the same value at $q_1$ and $q_2$ by the choice of $c$.
  \end{proof}

  Equality holds in~(ii) for every set of admissible pairs with
  $\Delta \le 6$ (by exhaustive enumeration), hence for all
  molecular graphs. It fails from $\Delta = 7$. The single violating set
  of at most four pairs there is $\{(1,7),(2,6),(3,5),(4,4)\}$, whose
  degree sums all equal $8$. A constant degree sum makes $\ln(i+j)$
  constant on the set, so $A \le 2$, while the bound equals $3$ as soon as
  the set has at least four pairs and three imbalances. A connected graph
  all of whose edges have the same degree sum $s$ realises only one pair.
  Indeed, each neighbour of a vertex of degree $d$ then has degree $s - d$, so by
  connectivity only the pair $\{d, s-d\}$ occurs. Hence no connected
  graph has this set as its pair set.
  As an illustration of~(iii), the spider with four legs of length three realises the
  pairs $(2,4)$, $(2,2)$ and $(1,2)$, which have only two imbalances
  ($\tfrac13$ and $0$). On it, $\gamma$ is confounded with
  $(\alpha,\beta)$ along the direction of part~(iii).
  For a QSPR dataset the relevant set is the union of the pairs realised
  by its molecules. A direction along which every molecule's index
  is rescaled by the same factor is invisible to any correlation or
  regression on the dataset. This concerns identifiability of $\theta$
  from the index values on the given set of graphs rather than
  identifiability of a fitted regression model. On the octanes this union contains nine of the ten
  admissible pairs and its affine dimension is $3$, so no such direction
  exists.

  The behaviour at the two ends of the $\gamma$ axis is used in the
  proof of Proposition~\ref{prop:laguerre}. Routine termwise bounds in
  $(m,\Delta,\delta)$, a Cauchy--Schwarz interpolation bound and a
  uniform ratio bound are collected in Supplementary Section~S9.

  \begin{proposition}\label{prop:gamma_eq_0_or_inf}
  At $\gamma = 0$, $LO(G; \alpha, \beta, 0)$ coincides with the
  two-parameter $(\alpha, \beta)$ degree-based index
  $\sum_{uv \in E(G)} (d_u d_v)^{\alpha}(d_u + d_v)^{\beta}$.
  Let $E_{\max}(G)$ denote the set of edges realising
  $\mathrm{imb}_{\max}(G) = \max_{uv} \mathrm{imb}(d_u, d_v)$
  and $E_{\min}(G)$ the set realising
  $\mathrm{imb}_{\min}(G) = \min_{uv} \mathrm{imb}(d_u, d_v)$.
  Then as $\gamma \to +\infty$,
  \[
  \frac{\Wuzi}{e^{\gamma\, \mathrm{imb}_{\max}(G)}}
  \;\longrightarrow\;
  \sum_{uv \in E_{\max}(G)} (d_u d_v)^{\alpha}(d_u + d_v)^{\beta},
  \]
  and as $\gamma \to -\infty$,
  \[
  \Wuzi / e^{\gamma\, \mathrm{imb}_{\min}(G)} \longrightarrow
  \sum_{uv \in E_{\min}(G)} (d_u d_v)^{\alpha}(d_u + d_v)^{\beta}.
  \]
  \end{proposition}

  \begin{proof}
  At $\gamma = 0$ the exponential factor is $1$ on every edge. With
  $\mu(G) := \mathrm{imb}_{\max}(G)$ and
  $\varepsilon_{uv} := \mu(G) - \mathrm{imb}(d_u, d_v) \ge 0$,
  we have $\Wuzi / e^{\gamma \mu(G)} = \sum_{uv} (d_u d_v)^{\alpha}(d_u +
  d_v)^{\beta} e^{-\gamma \varepsilon_{uv}}$. The summands over
  $E_{\max}(G)$ ($\varepsilon_{uv} = 0$) do not depend on $\gamma$, while the
  others ($\varepsilon_{uv} > 0$) tend to $0$ as $\gamma \to
  +\infty$, termwise since $E(G)$ is finite. The case
  $\gamma \to -\infty$ is dual.
  \end{proof}

% =====================================================================
\section{Discrimination}\label{sec:discrimination}
% =====================================================================


\subsection{The degree-pair ceiling}\label{sec:ceiling}

No vertex-degree-based index separates graphs with equal degree-pair
counts, so the number of distinct profiles is a \emph{ceiling} on
distinct values and a \emph{floor} on degeneracy.

  \begin{proposition}[Degree-pair information ceiling]\label{prop:collision}
  If two graphs $G$ and $H$ have identical unordered
  edge-degree-pair count vectors, $m_{ij}(G) = m_{ij}(H)$ for all
  $i \le j$, then $LO(G;\theta) = LO(H;\theta)$ for every
  $\theta \in \mathbb{R}^3$, and more generally
  $TI(G) = TI(H)$ for every vertex-degree-based index $TI$. Among the
  $18$ octane isomers exactly $16$ distinct count vectors occur, and the
  pairs $\{$3-methylheptane, 4-methylheptane$\}$ and
  $\{$3,4-dimethylhexane, 3-ethyl-2-methylpentane$\}$ share their
  count vectors but differ in measured properties
  (Table~\ref{tab:collisions}). Hence every vertex-degree-based
  descriptor, tuned or not, carries a nonzero irreducible error on
  this benchmark.
  \end{proposition}

  \begin{proof}
  By Definition~\ref{def:wuzi}, $LO(G;\theta)$ depends on $G$ only
  through the multiset of unordered edge degree pairs, that is, through
  $(m_{ij}(G))_{i \le j}$, and the same holds for every vertex-degree-based
  index~\cite{Rada2026LinearVDB}. The counts for the $18$ octane
  isomers are computed and cross-checked in the reproducibility
  package, and Table~\ref{tab:collisions} lists the common count
  vectors of both collision pairs.
  \end{proof}

  The ceiling is tight. Distinct count vectors are separated at almost
  every parameter choice.

  \begin{proposition}[Generic discrimination of distinct
  profiles]\label{prop:generic}
  Fix $\Delta \ge 1$ and let $m \neq m'$ be two distinct
  edge-degree-pair count vectors supported on the admissible pairs
  $P_{\Delta}$. Then
  $D(\theta) = \sum_{(i,j) \in P_{\Delta}} (m_{ij} - m'_{ij})\,
  e^{\theta^{\top} z_{ij}}$ is a real-analytic function of $\theta$
  that is not identically zero, and its zero set has Lebesgue measure
  zero in $\mathbb{R}^3$. Consequently, any two graphs with distinct
  count vectors are distinguished by $LO(\cdot\,;\theta)$ for almost
  every $\theta \in \mathbb{R}^3$, while only graphs with identical count
  vectors (Proposition~\ref{prop:collision}) are indistinguishable
  for all $\theta$.
  \end{proposition}

  \begin{proof}
  The map $(i,j) \mapsto z_{ij}$ is injective on $P_{\Delta}$, since
  $\ln(ij)$ and $\ln(i+j)$ recover the product and the sum, which
  determine the unordered pair. Choose $v \in \mathbb{R}^3$ such that
  the inner products $v^{\top} z_{ij}$ are pairwise distinct, which
  excludes only finitely many hyperplanes of $v$. Along the ray
  $\theta = t v$, $D(tv) = \sum_k c_k e^{t \lambda_k}$ with distinct
  exponents $\lambda_k$ and coefficients
  $c_k = m_{i_kj_k} - m'_{i_kj_k}$ that are not all zero, and such an
  exponential sum vanishes identically only if all $c_k = 0$. Hence $D \not\equiv 0$
  on the ray, and so on $\mathbb{R}^3$. Since $D$ is a finite sum of
  exponentials of linear functionals, it is real-analytic on
  $\mathbb{R}^3$, and the zero set of a nonzero real-analytic
  function has Lebesgue measure zero.
  \end{proof}

  The proof uses only the injectivity of
  $(i,j) \mapsto (\ln(ij), \ln(i+j))$ and therefore applies to the parent
  $M_{\alpha,\beta}$ on the $(\alpha,\beta)$-plane, so
  Proposition~\ref{prop:generic} is inherited from the parent rather
  than conferred by $\gamma$. Since $q_{ij} = \sqrt{1 - 4ij/(i+j)^{2}}$
  depends only on the pair, $\gamma$ adds no graph information beyond
  the edge-degree-pair counts, of which every vertex-degree-based index
  is a functional~\cite{Rada2026LinearVDB}. What it
  adds is the functional direction of Lemma~\ref{prop:indep}, that is,
  a different weighting of the same information. As a corollary, on graphs of maximum degree four any
  ten degree-based indices with linearly independent weight vectors
  $(f(i,j))_{i \le j \le 4}$ span the ten-dimensional count
  space. Hence every bond-incident-degree index, the Loyola family at
  every parameter value included, is an exact linear combination of
  them and adds no linear information for any target. Whether such a reweighting still helps prediction is the question of
  Section~\ref{sec:sensitivity}.

\begin{table}[htbp]
\centering
\caption{The two degree-pair collision pairs among the $18$ octane
isomers. Each pair shares its complete edge-degree-pair count vector,
shown as multiplicities of the unordered pairs, and is therefore
indistinguishable by every vertex-degree-based index for all
parameter values. The two members differ in the measured properties.}
\label{tab:collisions}
\small
\begin{tabular}{lrrrrr}
\toprule
Isomer & $T_B$ & $\Delta H_f$ & $\Delta H_{\mathrm{vap}}$ & $S$ & $\omega$\\
\midrule
\multicolumn{6}{l}{\emph{shared count vector} $2(1,2),\,1(1,3),\,2(2,2),\,2(2,3)$:}\\
3-methylheptane & $118.9$ & $-50.82$ & $9.52$ & $111.26$ & $0.371$\\
4-methylheptane & $117.7$ & $-50.69$ & $9.48$ & $109.32$ & $0.372$\\
\midrule
\multicolumn{6}{l}{\emph{shared count vector} $2(1,2),\,2(1,3),\,2(2,3),\,1(3,3)$:}\\
3,4-dimethylhexane & $117.7$ & $-50.91$ & $9.32$ & $106.59$ & $0.340$\\
3-ethyl-2-methylpentane & $115.6$ & $-50.48$ & $9.21$ & $106.06$ & $0.330$\\
\bottomrule
\end{tabular}
\end{table}

\subsection{Degeneracy across tree orders}\label{sec:degeneracy}

Degeneracy~\cite{Konstantinova1996Degeneracy} measures the failure to
discriminate non-isomorphic graphs, and we quantify it with the
distinct-value convention $100(1 - k/N)$, where $k$ is the number of
distinct values among $N = 106$ trees (Supplementary Section~S5). On
the $106$ trees of order $10$ the pure-$\gamma$ points $LO(0, 0, 1)$
and $LO(0, 0, 2)$ attain $25.5\%$ ($79$ distinct values), tied with
$R$, $GA$ and $AG$. The $106$ trees realise only $79$ distinct
edge-degree-pair count vectors, so $25.5\%$ is the floor for
\emph{any} degree-pair-based index, while $M_1$ ($83.0\%$) and
$M_2$ ($67.0\%$) are the most degenerate (Supplementary Figure~S2).
For orders $7$ through $12$, on all trees and on those with
$\Delta \le 4$ (Supplementary Table~S4), the pure-$\gamma$ points and
$R$ attain the floor at every order. The index $\chi$ leaves it from
order $10$, with $77$ values against $79$ profiles.
Theorem~\ref{thm:ratioplane} makes this exact. The pure-$\gamma$
points separate precisely the $q$-histograms, which coincide with the
degree-pair profiles up to order $12$ (order $15$ for $\Delta \le 4$)
and not beyond. At order $13$ the profiles
$\{(1,2)^{1}, (1,3)^{2}, (1,6)^{5}, (2,3)^{2}, (2,6)^{1}, (3,3)^{1}\}$
and $\{(1,3)^{3}, (1,6)^{5}, (2,2)^{1}, (2,3)^{2}, (3,6)^{1}\}$
(exponents are edge multiplicities) share the $q$-histogram
$\{0^{1}, \tfrac15^{2}, \tfrac13^{1}, \tfrac12^{3}, \tfrac57^{5}\}$.
Table~\ref{tab:firstfail} collects the first failures, and the index
$R$ still attains the floor at order $13$ ($570$ values) and leaves it
at order $14$, where the pure-$\gamma$ points give $1090$ values.

\begin{table}[htbp]
\centering
\caption{First tree order at which discrimination is lost. Values are
counted as distinct at $11$ significant digits. ``Below
floor'' is the first order with fewer distinct values than distinct
edge-degree-pair profiles, the floor for every vertex-degree-based
index, and the entry gives (values/profiles) at that order. ``Merges'' is
the first order at which an index on $\Pi$ gives equal values to two
trees with different $q$-histograms. This does not apply to $R$, which
is not on $\Pi$. $LO(0,0,\gamma)$ denotes the pure-$\gamma$ points
$\gamma = 1, 2$. Orders below $7$ were not examined. A dash entry means
not determined.}
\label{tab:firstfail}
\small
\setlength{\tabcolsep}{3pt}
\begin{tabular}{lcccc}
\toprule
 & \multicolumn{2}{c}{all trees} & \multicolumn{2}{c}{trees with $\Delta \le 4$}\\
\cmidrule(lr){2-3}\cmidrule(lr){4-5}
Index & below floor & merges & below floor & merges\\
\midrule
$LO(0,0,\gamma)$ & $13$ ($566/570$) & never$^{a}$ & $16$ ($809/810$) & never$^{a}$\\
$R$ & $14$ ($1095/1100$) & n/a$^{b}$ & $> 16^{c}$ & n/a$^{b}$\\
$GA$ & $13$ ($566/570$)$^{d}$ & $17^{e}$ & $16$ ($809/810$)$^{d}$ & ---$^{e}$\\
\bottomrule
\end{tabular}

\smallskip
\begin{minipage}{0.95\linewidth}\footnotesize
$^{a}$~All $q$-histograms are separated for every nonzero algebraic
$\gamma$ (Theorem~\ref{thm:ratioplane}(ii)). $^{b}$~$R$ is not on $\Pi$.
$^{c}$~$R$ attains the floor for trees with $\Delta \le 4$ at every
order up to $16$, the largest examined. $^{d}$~$GA$ lies on $\Pi$
and merges no $q$-histograms below order $17$, so there it ties the
pure-$\gamma$ points. $^{e}$~The witness (Section~\ref{sec:ratioplane})
has maximum degree $5$.
\end{minipage}
\end{table}

Off the ratio plane the imbalance factor can restore discrimination
lost at particular classical exponents. A point $(0,\beta,0)$ weights
an edge by its degree sum only, so pairs with equal sum receive equal
weight. With $\gamma \neq 0$ their imbalances separate them. On all
trees of order $16$ ($4069$ profiles), $LO(0,\tfrac12,0)$ and
$LO(0,\tfrac13,0)$ take only $3475$ and $3476$ distinct values, while
$LO(0,\tfrac12,1)$ and $LO(0,\tfrac13,1)$ separate all $4069$.
$LO(\tfrac12,0,\cdot)$ rises from $3940$ to $4067$ and the Randi\'c
point $LO(-\tfrac12,0,\cdot)$ from $4007$ to $4068$ as $\gamma$ goes
from $0$ to $1$. Points on $\Pi$ such as $GA$ are unaffected ($3996$
in both cases), as Theorem~\ref{thm:ratioplane} requires, and generic
parent points such as $LO(0.3,0.2,0)$ already separate all $4069$
profiles (Proposition~\ref{prop:generic}). The effect is established
by enumeration for orders $12$--$16$ but is not proved. It gives
no predictive gain over a tuned parent in Section~\ref{sec:sensitivity}.
At fixed fractional exponents the sign of the change depends on the
property (Supplementary Section~S5).

\subsection{Structure sensitivity}\label{sec:ss}

The structure-sensitivity criterion of Furtula, Gutman and
Dehmer~\cite{FurtulaGutmanDehmer2013SS}, with a fingerprint-based
variant due to Raki\'c and Furtula~\cite{RakicFurtula2019SS}, was
implemented with the published protocol on the $75$ decane isomer
trees. For a tree, the similar structures are those obtained by
relocating a single edge within the molecular class. $SS$ and $Abr$
are the class means of the mean and the maximum relative change of the
index over these neighbours. The implementation reproduces the
published decane-class values of Barman and Das~\cite{BarmanDas2026HSO},
with eleven of fourteen $SS$/$Abr$ values agreeing exactly at the four
displayed decimals. Each of the other three agrees within one unit in
the fourth. The protocol and this control are detailed in Supplementary
Section~S6.
Under the published desideratum ($SS$ large, $Abr$ small), $LO(0,0,2)$
has a structure sensitivity indistinguishable at three decimals from
that of the hyper-Zagreb index
($0.1257$ against $0.1256$). Its abruptness is markedly smaller ($0.274$
against $0.313$). The sensitivity-to-abruptness ratios of the
pure-$\gamma$ points ($0.4579$ and $0.4565$) are comparable to those of
$H/2$, $GA$ and $R$ ($0.448$--$0.451$, Table~\ref{tab:fgdss}). In view
of the redundancy noted in Section~\ref{sec:octanes}, these
third-decimal differences do not rank the indices. $SS$ is also not a
fixed property of the family. Along the pure-$\gamma$ line it grows
with $\gamma$ ($0.031$, $0.062$, $0.126$, $0.319$ at
$\gamma = 0.5, 1, 2, 5$), while the ratio peaks near $\gamma = 2$ and
then falls ($0.427$ at $\gamma = 5$). The pure-$\gamma$ points reduce
to first order in $\gamma$ to the parameter-free aggregate
$IRLA$, since $LO(0,0,\gamma) = m + \tfrac{\gamma}{2}\,IRLA + O(\gamma^2)$, and this aggregate has a higher ratio ($SS = 0.2343$, $Abr = 0.5066$, ratio
$0.4625$). Hence the favourable sensitivity--abruptness balance is
a property of linear imbalance weighting rather than of the
exponential form. The equality of $SS$ with hyper-Zagreb at
$\gamma = 2$ reflects the chosen value of $\gamma$.

\begin{table}[htbp]
\centering
\caption{Structure sensitivity $SS$, abruptness $Abr$, and their
ratio on the $75$ decane isomers under the published
similar-structure protocol of~\cite{FurtulaGutmanDehmer2013SS}
(single-edge relocation within the molecular class).}
\label{tab:fgdss}
\small
\begin{tabular}{lrrr}
\toprule
Index & $SS$ & $Abr$ & $SS/Abr$\\
\midrule
$M_1$ & $0.0566$ & $0.1336$ & $0.4238$\\
$M_2$ & $0.0922$ & $0.2443$ & $0.3776$\\
$HM$ & $0.1256$ & $0.3126$ & $0.4017$\\
${}^{m}\!M_2$ & $0.0516$ & $0.1177$ & $0.4384$\\
$R$ & $0.0264$ & $0.0588$ & $0.4481$\\
$\chi$ & $0.0264$ & $0.0591$ & $0.4463$\\
$H/2$ & $0.0516$ & $0.1143$ & $0.4514$\\
$ISI$ & $0.0226$ & $0.0597$ & $0.3795$\\
$GA$ & $0.0241$ & $0.0535$ & $0.4507$\\
$AG$ & $0.0256$ & $0.0579$ & $0.4415$\\
$LO(0,0,1)$ & $0.0624$ & $0.1367$ & $0.4565$\\
$LO(0,0,2)$ & $0.1257$ & $0.2745$ & $0.4579$\\
\bottomrule
\end{tabular}
\end{table}

% =====================================================================
\section{QSPR diagnostics}\label{sec:sensitivity}
% =====================================================================

In Tables~\ref{tab:box}--\ref{tab:bp},
medians are taken over candidate-stream seeds, and seed counts summarise the
variability of the random search on fixed data and are not replications.
Values are computed in double precision and rounded once for display
(three decimals for correlations, $Q^2$ and RMSE, four for structure
sensitivity). A tabulated difference may therefore differ in the last place
from the difference of two displayed entries, and a quoted median of
differences need not equal the difference of medians. The narrow- and
wide-box streams use the same uniform variates scaled by the half-width,
so comparisons across boxes are paired.

We benchmark the Loyola family against the classical indices it
reduces to on the $18$ octane isomers, following the chemometric
benchmark of the diminished Sombor
study~\cite{MovahediGutmanRedzepovicFurtula2026}. The $35$
nonane isomers are a replication set, with models refitted on it. The
comparison set consists of the exact reductions of Section~\ref{sec:intro},
$M_1, M_2, HM, {}^{m}\!M_2, R, \chi, H/2, ISI, GA, AG$, and the two
fixed pure-$\gamma$ benchmark points $LO(0, 0, 1)$ and $LO(0, 0, 2)$.
Since the edge count $m = LO(0, 0, 0)$ is constant on the octanes, it is
omitted. No parameter in the correlation comparison is tuned against
the property data. An exploratory local-tuning check (Supplementary
Section~S2) did not exceed the best fixed benchmark except by a $0.001$ in-sample
$\Delta H_{\mathrm{vap}}$ margin. All comparisons are
descriptive. On sets this small they indicate how imbalance information
changes the fit of the parent family, and they do not establish
chemical predictive utility.

\subsection{Correlations on the octane isomers}\label{sec:octanes}

Table~\ref{tab:lo_octane} reports the Pearson correlation $r$ between
each index and each of the five physicochemical properties of the
$18$ octane isomers. The properties are the boiling point $T_B$, the enthalpy of formation
$\Delta H_f$, the enthalpy of vaporization $\Delta H_{\mathrm{vap}}$, the
entropy $S$ and the acentric factor $\omega$. Data and units are in
Table~\ref{tab:octane-data}. Boiling points, formation enthalpies and
(with one exception) vaporization enthalpies are traced to the NIST
WebBook~\cite{NISTWebBook}, and acentric factors to KDB-derived
compilations. The WebBook lists a gas-phase entropy for only five
isomers, while seven entropies are compared with the 1947 API Research
Project~44 tables~\cite{APIProject44Circ461} or a compiled octane
dataset. All $90$ molecule-level values were audited against
authoritative sources (Supplementary Section~S1).
Three caveats remain. The $298$\,K vaporization enthalpy of
2,2,3,3-tetramethylbutane, a solid at $298$\,K (triple point
$374$\,K), is the benchmark-compilation value $8.41$\,kcal\,mol$^{-1}$
($35.19$\,kJ\,mol$^{-1}$). This matches neither of the NIST standard vaporization
enthalpies ($42.94$ and $42.91$\,kJ\,mol$^{-1}$, about
$10.26$\,kcal\,mol$^{-1}$), so its physical reference state is
unresolved. Six gas-phase entropies could not be verified against a
primary source, and the 1947 API Project 44 values differ from them by up to
$2.2$ units. The 2,3,3-trimethylpentane value is a compiled-dataset
value, and four entropies differ from a widely used compilation.
Substituting the compilation values changes all twelve $S$ correlations (by up to $0.043$ in $|r|$)
and moves the margin between $HM$ and $M_1$ at the top from $0.001$ to
$0.007$. Margins of this size are not meaningful at this resolution. Alternative compiled
acentric factors change several $\omega$ correlations in the third
decimal (e.g.\ $M_2$ from $-0.988$ to $-0.986$) and leave the $\omega$
ranking unchanged (detail in Supplementary Section~S10).

Table~\ref{tab:octane-data} lists the full $18$-isomer dataset used
in this section, with $T_B$ in $^{\circ}$C, $\Delta H_f$ and
$\Delta H_{\mathrm{vap}}$ in kcal\,mol$^{-1}$
($\Delta H_{\mathrm{vap}}$ at $298$~K), $S$ in
cal\,mol$^{-1}$K$^{-1}$ and $\omega$ dimensionless. The
$\Delta H_{\mathrm{vap}}$ values of 3,3-dimethylhexane ($8.97$),
3-ethyl-3-methylpentane ($9.08$), 2,2,3-trimethylpentane ($8.83$) and
2,3,3-trimethylpentane ($8.90$) are values consistent with the NIST
WebBook determinations, and Supplementary Section~S1 gives the
exact-versus-within-uncertainty status of each. The
$\Delta H_{\mathrm{vap}}$ entry of 2,2,3,3-tetramethylbutane is the
benchmark-compilation value discussed above. The $S$ value of
2,3,3-trimethylpentane ($102.06$) is a compiled-dataset value
(Supplementary Section~S10).


Table~\ref{tab:octane-data} has several open provenance items, namely
the unverified entropies, the vaporization enthalpy of
2,2,3,3-tetramethylbutane, the corrected transcription of the
2,3,3-trimethylpentane entropy and the alternative compiled values.
These items, with their effect on Table~\ref{tab:lo_octane}, are
collected in Supplementary Section~S10, and none of them enters the
pre-registered tests of Sections~\ref{sec:external} and~\ref{sec:bp}.

\begin{table}[H]
\centering\small
\caption{Physicochemical properties of the $18$ octane isomers.}
\label{tab:octane-data}
\begin{tabular}{lrrrrr}
\toprule
Isomer & $T_B$ & $\Delta H_f$ & $\Delta H_{\mathrm{vap}}$ & $S$ & $\omega$\\
\midrule
$n$-octane & 125.6 & $-49.82$ & 9.92 & 111.55 & 0.398\\
2-methylheptane & 117.6 & $-51.50$ & 9.48 & 109.84 & 0.378\\
3-methylheptane & 118.9 & $-50.82$ & 9.52 & 111.26 & 0.371\\
4-methylheptane & 117.7 & $-50.69$ & 9.48 & 109.32 & 0.372\\
2,2-dimethylhexane & 106.8 & $-53.71$ & 8.92 & 103.13 & 0.339\\
2,3-dimethylhexane & 115.6 & $-51.13$ & 9.27 & 108.02 & 0.348\\
2,4-dimethylhexane & 109.4 & $-52.44$ & 9.03 & 106.98 & 0.344\\
2,5-dimethylhexane & 109.1 & $-53.21$ & 9.05 & 105.72 & 0.357\\
3,3-dimethylhexane & 111.9 & $-52.61$ & 8.97 & 104.74 & 0.322\\
3,4-dimethylhexane & 117.7 & $-50.91$ & 9.32 & 106.59 & 0.340\\
3-ethyl-2-methylpentane & 115.6 & $-50.48$ & 9.21 & 106.06 & 0.330\\
3-ethyl-3-methylpentane & 118.3 & $-51.38$ & 9.08 & 101.48 & 0.302\\
3-ethylhexane & 118.5 & $-50.40$ & 9.48 & 109.43 & 0.362\\
2,2,3-trimethylpentane & 109.8 & $-52.61$ & 8.83 & 101.31 & 0.300\\
2,2,4-trimethylpentane & 99.2 & $-53.57$ & 8.40 & 101.81 & 0.305\\
2,3,3-trimethylpentane & 114.8 & $-51.73$ & 8.90 & 102.06 & 0.291\\
2,3,4-trimethylpentane & 113.5 & $-51.97$ & 9.01 & 102.39 & 0.317\\
2,2,3,3-tetramethylbutane & 106.5 & $-53.99$ & 8.41 & 93.06 & 0.247\\
\bottomrule
\end{tabular}
\end{table}



% References should be ordered alphabetically.

\begin{table}[htbp]
\centering
\caption{Pearson correlation $r$ between the Loyola reductions, the fixed pure-$\gamma$ benchmark points and the five physicochemical properties of the $18$ octane isomers. Bold flags the largest $|r|$ per column, and this is descriptive. Margins to the runner-up are $0.001$--$0.032$, and their paired bootstrap intervals include zero (see the text).}
\label{tab:lo_octane}
\small
\begin{tabular}{lrrrrr}
\toprule
Index & $T_B$ & $\Delta H_f$ & $\Delta H_{\mathrm{vap}}$ & $S$ & $\omega$\\
\midrule
$M_1 = LO(0,1,0)$ & $-0.718$ & $-0.762$ & $-0.935$ & $-0.970$ & $-0.970$\\
$M_2 = LO(1,0,0)$ & $-0.497$ & $-0.542$ & $-0.811$ & $-0.934$ & $\mathbf{-0.988}$\\
$HM = LO(0,2,0)$ & $-0.654$ & $-0.710$ & $-0.903$ & $\mathbf{-0.971}$ & $-0.982$\\
${}^{m}\!M_2 = LO(-1,0,0)$ & $\mathbf{+0.855}$ & $\mathbf{+0.891}$ & $+0.925$ & $+0.863$ & $+0.794$\\
$R = LO(-1/2,0,0)$ & $+0.819$ & $+0.850$ & $+0.958$ & $+0.932$ & $+0.898$\\
$\chi = LO(0,-1/2,0)$ & $+0.800$ & $+0.832$ & $+0.961$ & $+0.947$ & $+0.924$\\
$H/2 = LO(0,-1,0)$ & $+0.821$ & $+0.849$ & $+0.961$ & $+0.932$ & $+0.900$\\
$ISI = LO(1,-1,0)$ & $-0.005$ & $-0.000$ & $-0.383$ & $-0.593$ & $-0.738$\\
$GA = 2\,LO(1/2,-1,0)$ & $+0.822$ & $+0.858$ & $+0.965$ & $+0.939$ & $+0.906$\\
$AG = \tfrac12 LO(-1/2,1,0)$ & $-0.816$ & $-0.859$ & $-0.957$ & $-0.938$ & $-0.900$\\
$LO(0,0,1)$ & $-0.828$ & $-0.832$ & $\mathbf{-0.984}$ & $-0.928$ & $-0.929$\\
$LO(0,0,2)$ & $-0.829$ & $-0.846$ & $-0.981$ & $-0.937$ & $-0.925$\\
\bottomrule
\end{tabular}
\end{table}

The best correlation per property is split, with ${}^{m}\!M_2 = LO(-1, 0, 0)$
strongest on $T_B$ and $\Delta H_f$, $HM = LO(0, 2, 0)$ on $S$ and
$M_2 = LO(1, 0, 0)$ on $\omega$. The fixed pure-$\gamma$ point
$LO(0, 0, 1)$ is strongest on $\Delta H_{\mathrm{vap}}$ ($|r| = 0.984$, with
$LO(0, 0, 2)$ second at $0.981$), and at the exact-reduction triples the
Loyola correlations coincide with the classical indices, as they must.
For each property a descriptive post-selection paired bootstrap
interval for the observed top-two $|r|$ difference included zero, but these
intervals are not multiplicity-adjusted and do not establish equivalence.
Recomputing each bolded correlation with single isomers deleted
shifts $|r|$ by at most $0.027$ (bootstrap details are in Supplementary
Section~S1). The bolded entries are the largest of twelve highly correlated sample
correlations on $18$ molecules, and are therefore optimistic (winner's-curse)
estimates of the selected descriptor's correlation on new molecules.
They are descriptive rankings rather than tests.

Fourteen of the descriptor pairs
have $|r| \ge 0.99$, and the first two principal components carry
$99.05\%$ of the variance (Supplementary Section~S1). The
pure-$\gamma$ points are collinear with the classical reductions
(maximum $|r|$ with a reduction equal to
$0.982$ for $LO(0, 0, 1)$ and $0.993$ for $LO(0, 0, 2)$). Differences
in the third decimal of $|r|$ between near-collinear descriptors are
therefore not meaningful rankings.

Proposition~\ref{prop:collision} is realised on these data, where
the $18$ isomers realise only $16$ distinct edge-degree-pair count
vectors. The two collision pairs of Table~\ref{tab:collisions}
are indistinguishable by \emph{every} vertex-degree-based index at
\emph{every} parameter choice. Their measured properties
differ, by up to $1.2\,^{\circ}$C in $T_B$ within the first pair and
$2.1\,^{\circ}$C within the second. Every correlation in
Table~\ref{tab:lo_octane}, and every tuned variant below, operates
above this irreducible floor.

\subsection{Effect of freeing $\gamma$ and of the search
box}\label{sec:ablation}

We ask whether, at identical search budgets, freeing $\gamma$ improves
out-of-sample prediction beyond an optimised parent $M_{\alpha,\beta}$.
The design is a fully nested leave-one-out (details in Supplementary
Section~S3) in which, for each property and each outer fold, candidate parameter
vectors are drawn uniformly per coordinate from a search box. These are triples
$(\alpha,\beta,\gamma)$ for the Loyola model and pairs $(\alpha,\beta)$
with $\gamma = 0$ for the parent, with identical budget. Candidates are scored on the
training isomers only, by \emph{inner} leave-one-out root-mean-square
error of the one-descriptor linear regression, after which the best candidate is
refit and the held-out isomer predicted. No quantity computed on the
held-out isomer enters selection or fitting. In the box $(-2,2)$ used
first, freeing $\gamma$ appeared to help the vaporization enthalpy
(Supplementary Table~S3), improving
nested predictive performance for $\Delta H_{\mathrm{vap}}$ in every
seed over one hundred candidate streams. The count was $100/100$ at budgets $200$ and $500$, with median
$\Delta Q^2 = +0.069$ and $+0.068$, whereas freeing $\gamma$ usually degraded
$\Delta H_f$ and gave mixed or slightly negative results for the other
properties. With 2,2,3,3-tetramethylbutane excluded the count falls to
$93/100$ and $96/100$ (Supplementary Section~S3), and these counts describe
optimiser variance on one fixed set of $18$ molecules. Selection also varies across folds. For $\Delta H_f$
the three-parameter search chose three distinct triples over the $18$
folds, whereas for $\Delta H_{\mathrm{vap}}$ it chose the same triple
in every fold (Supplementary Section~S3). The nested $Q^2$ therefore evaluates
the selection-plus-fit procedure rather than one fitted model. These counts also
hide a boundary effect, since for
$\Delta H_{\mathrm{vap}}$ the parent chose
$(\alpha,\beta) = (1.021, -1.993)$, on the boundary of the box, in $16$
of the $18$ folds.

The decisive design choice is the box $(-2,2)$, which the nonane
pre-registration also fixed. By $\ln(ij) = 2\ln(i+j) +
\ln(1-q^2) - \ln 4$ (Eq.~\eqref{eq:factor}) the parent family
reaches imbalance information through $(1-q^2)^{\alpha}$ along the
direction $\beta \approx -2\alpha$. It needs large exponents to
weight it strongly, and the box truncates exactly that direction. The
matched comparison was therefore repeated with $2000$
candidates per model, seeds $0$--$49$, and boxes of half-width $2$
and $12$ for both models (Table~\ref{tab:box}), a repetition that was
not pre-registered. In the box of
half-width $12$ the parent's selections are interior in more than
$90\%$ of folds for every property, and the vaporization-enthalpy gain
shrinks to about zero. A hybrid control that widens only
$(\alpha,\beta)$ gives median $\Delta Q^2 = +0.004$ on the octanes
($28$ of $50$ seeds) and $-0.002$ on the nonanes ($24$ of $50$).
Widening all three coordinates gives $-0.010$ and $-0.016$, while on the
other properties no advantage is consistent across boxes. Under the
hybrid control the Loyola model is ahead for octane $S$ ($38$ of $50$
seeds) and nonane $T_B$ ($31$ of $50$), and behind for octane $T_B$,
$\Delta H_f$ and $\omega$. The wide box itself lowers the parent's
$Q^2$ for $S$ and $T_B$. The vaporization-enthalpy gains above are thus
a truncation effect rather than information the parent lacks. Most of
the change comes from the parent improving in
the wide box ($0.887 \to 0.942$ on the octanes). The rest comes from the
three-dimensional search being slightly less efficient there ($0.957
\to 0.946$ for the hybrid control). What remains is due to linear imbalance weighting itself. Both the fixed point $LO(0,0,1)$ and the
parameter-free $IRLA$ approach what the parent reaches only near
$(\alpha,\beta) \approx (4.5, -9)$ ($Q^2 = 0.957$ and $0.948$ against
$0.949$ there, on the octanes).

\begin{table}[htbp]
\centering
\caption{Sensitivity of the matched comparison to the search box.
Oct.: octanes, $n = 18$. Non.: nonanes. Edge: share of outer folds in
which the parent's selection lies within $10\%$ of the boundary of the
box of the stated half-width. Medians of $Q^2$ are over seeds $0$--$49$ at $2000$
candidates per model. Counts: the number of seeds (of $50$) in which LO
beats GM (half-width $2$) or LO$_h$ beats GM (half-width $12$). LO$_h$ is
a hybrid control with $(\alpha,\beta)$ in the wide box and $\gamma$ in
$(-2,2)$, which avoids diluting the Loyola search over a wide $\gamma$
range.}
\label{tab:box}
\small
\setlength{\tabcolsep}{1.6pt}
\begin{tabular}{lrrrrrrrrr}
\toprule
 & \multicolumn{4}{c}{half-width $2$} & \multicolumn{5}{c}{half-width $12$}\\
\cmidrule(lr){2-5}\cmidrule(lr){6-10}
Set, prop. & Edge & GM & LO & $>$ & Edge & GM & LO & LO$_h$ & $>$\\
\midrule
oct. $T_B$ & $69\%$ & $0.665$ & $0.730$ & $36$ & $2\%$ & $0.606$ & $0.495$ & $0.544$ & $21$\\
oct. $\Delta H_f$ & $52\%$ & $0.821$ & $0.769$ & $11$ & $0\%$ & $0.820$ & $0.670$ & $0.761$ & $10$\\
oct. $\Delta H_{\mathrm{vap}}$ & $97\%$ & $0.887$ & $0.957$ & $50$ & $3\%$ & $0.942$ & $0.934$ & $0.946$ & $28$\\
oct. $S$ & $85\%$ & $0.904$ & $0.891$ & $11$ & $4\%$ & $0.842$ & $0.838$ & $0.867$ & $38$\\
oct. $\omega$ & $0\%$ & $0.979$ & $0.974$ & $10$ & $0\%$ & $0.979$ & $0.938$ & $0.951$ & $7$\\
non. $T_B$ & $61\%$ & $0.576$ & $0.601$ & $30$ & $0\%$ & $0.524$ & $0.392$ & $0.577$ & $31$\\
non. $\Delta H_{\mathrm{vap}}$ & $93\%$ & $0.855$ & $0.895$ & $49$ & $9\%$ & $0.892$ & $0.881$ & $0.889$ & $24$\\
\bottomrule
\end{tabular}
\end{table}

Under these protocols, then, tuning $\gamma$ yields no demonstrated
predictive advantage over a parent evaluated over the wider search box.
Equal candidate budgets do not guarantee equally thorough search in two
and three dimensions, and the hybrid control, which keeps $\gamma$ in a
narrow range, reduces this asymmetry without removing it. A more
exhaustive three-dimensional search could in principle find a better
Loyola point. The value of $\gamma$ lies in the structural results of
Section~\ref{sec:geometry} and in the comparison with linear imbalance
weighting of Section~\ref{sec:external}.

\subsection{Replication on the nonane isomers}\label{sec:external}

Under a protocol committed to the reproducibility package before any
nonane value was collected, the identical nested design was applied to
the $35$ constitutional nonane isomers. Boiling points were $35$ values:
$20$ NIST averages, $5$ single determinations and $10$ most recent
determinations selected under the protocol's Deviation D1. Standard vaporization
enthalpies were $34$ values from the NIST WebBook~\cite{NISTWebBook}. Only $4$ of the $34$ vaporization
enthalpies are calorimetric determinations (one more is a NIST average),
and $29$ come from two
compilations~\cite{LabaufGreenshieldsRossini1961, WilhoitZwolinski1971}.
Agreement with them therefore measures consistency with compiled values rather than
independent experimental validation. Within-fold baselines are in
Table~\ref{tab:external}. The pre-registered verdict is that the
$\gamma$ effect \emph{replicates} for $\Delta H_{\mathrm{vap}}$
(compilation values), with LO better than GM in $95/100$ and $98/100$ seeds
and median $\Delta Q^2 = +0.048$ and $+0.053$. For $T_B$ the verdict is \emph{mixed}
($58/100$ and $63/100$). The pre-registered protocol, however, fixed the
box of half-width $2$, which binds the parent (Table~\ref{tab:box}), and in
a box that does not, the nonane gain disappears (median
$\Delta Q^2 = -0.002$ under the hybrid control). The pre-registered
verdict is reported as obtained, with this deviation recorded. On both sets the
untuned fixed point $LO(0,0,1)$ matches or
beats the tuned Loyola model ($0.957$ and $0.903$ against
medians $0.953$ and $0.900$). It is closely related to $IRLA$:
$LO(G;0,0,\gamma) = m + \tfrac{\gamma}{2}\,IRLA(G) + O(\gamma^2)$ as $\gamma \to 0$.
At $\gamma = 1$ higher moments of the imbalance remain, and trees with
equal $IRLA$ can differ in $LO(0,0,1)$. Even so, the correlation between the two is $0.996$ on both sets. $IRLA$ itself
comes within $0.003$ of the median Loyola model on the nonanes ($0.897$
against $0.900$). What matters for this property is therefore a term linear in
the imbalance. Ridge regression on the full
degree-pair count vector predicts the nonanes better than every index
model ($Q^2 = 0.927$ for $\Delta H_{\mathrm{vap}}$, $0.682$ for $T_B$).
This is consistent with Proposition~\ref{prop:collision} and the remark after
Proposition~\ref{prop:generic} (every index in the family is a
one-dimensional reweighting of the degree-pair counts). A post hoc molecule
bootstrap, conditional on the fitted predictions, omits refitting
and selection uncertainty, and its $95\%$ interval for LO minus GM on nonane
$\Delta H_{\mathrm{vap}}$, $[-0.006, +0.113]$ at budget $200$, includes
zero. Protocol details, an exploratory transfer analysis and this
bootstrap are in Supplementary Section~S4.

\begin{table}[htbp]
\centering
\caption{Nested leave-one-out $Q^2$ on the octanes (oct., $n = 18$) and
the nonane replication set (non., $n = 35$ for $T_B$, $34$ for
$\Delta H_{\mathrm{vap}}$). Mean: training-fold mean. Class.: best of the
ten classical indices, selected inside each fold. IRLA: one-descriptor
model on $IRLA$. Ridge: ridge regression on the edge-degree-pair counts.
LO$_1$: the fixed, untuned point $LO(0,0,1)$. GM, LO: median of $Q^2$
over seeds $0$--$99$ at budget $200$ in the box of half-width $2$
(see Table~\ref{tab:box} for wider boxes). Last column: seeds in which
LO beat GM at budgets $200$/$500$.}
\label{tab:external}
\small
\setlength{\tabcolsep}{2pt}
\begin{tabular}{lrrrrrrrr}
\toprule
Set, property & Mean & Class. & IRLA & Ridge & LO$_1$ & GM & LO & LO$>$GM\\
\midrule
oct. $T_B$ & $-0.121$ & $0.593$ & $0.585$ & $0.344$ & $0.591$ & $0.662$ & $0.652$ & $53/64$\\
oct. $\Delta H_f$ & $-0.121$ & $0.746$ & $0.593$ & $0.654$ & $0.639$ & $0.831$ & $0.689$ & $14/24$\\
oct. $\Delta H_{\mathrm{vap}}$ & $-0.121$ & $0.884$ & $0.948$ & $0.954$ & $0.957$ & $0.888$ & $0.953$ & $100/100$\\
oct. $S$ & $-0.121$ & $0.905$ & $0.733$ & $0.721$ & $0.785$ & $0.904$ & $0.889$ & $34/31$\\
oct. $\omega$ & $-0.121$ & $0.971$ & $0.810$ & $0.851$ & $0.823$ & $0.979$ & $0.970$ & $26/28$\\
non. $T_B$ & $-0.060$ & $0.502$ & $0.400$ & $0.682$ & $0.409$ & $0.577$ & $0.586$ & $58/63$\\
non. $\Delta H_{\mathrm{vap}}$ & $-0.062$ & $0.857$ & $0.897$ & $0.927$ & $0.903$ & $0.853$ & $0.900$ & $95/98$\\
\bottomrule
\end{tabular}
\end{table}

\subsection{Measured boiling points}\label{sec:bp}

Measured vaporization enthalpies are too scarce for a larger
test. NIST lists only $25$ acyclic alkane skeletons outside $C_8$ and
$C_9$, and one of the $75$ decanes, with at least one calorimetric
standard-state $298$\,K determination (NIST averaged entries resolved
against their individual measurements). Two further tests therefore used measured boiling points, under a
protocol committed to the repository before any model was fitted. The
NIST pages had been archived two days earlier during a
data-availability audit that counted available values only. Boiling
points were taken from the NIST
WebBook~\cite{NISTWebBook} by a fixed rule. The rule takes the NIST average, or else
the median of non-compilation determinations if they agree within
$3$\,K, and otherwise excludes the molecule. This gives all hexanes, heptanes and octanes and $34$ of the $35$ nonanes,
since one nonane has two NIST entries with the same InChI whose boiling points differ
by $6.4$\,K. It also gives $34$ of the $75$ decanes, because $31$ of them have no NIST boiling point
and $10$ have determinations more than $3$\,K apart. For $14$ of the $34$ included decanes NIST lists a single
non-compilation determination, so the agreement rule does not apply to
them. Under this uniform rule the pooled octane boiling points can
differ from those of
Table~\ref{tab:octane-data} by up to $0.5\,^{\circ}$C, for example
2,4-dimethylhexane at $108.9$ against $109.4\,^{\circ}$C, a documented
source variation. Eight pooled nonane boiling points differ from
the replication-set values of Section~\ref{sec:external} by
$0.1$--$0.9\,^{\circ}$C, as that protocol took the most recent
determination where NIST gives no average. The included and excluded
decanes do not differ detectably in
branching or imbalance. The measures were the numbers of tertiary and quaternary carbons and
of methyl groups, the Wiener index, $IRLA$, $LO(0,0,1)$ and $R$. Two-sided
permutation tests gave $p \ge 0.10$ in every case, though with $34$ and $41$
molecules such tests can rule out only large differences. Test~A
evaluates the models on the decanes as a new cohort, refitting and
reselecting inside each leave-one-out fold. Test~B pools all
$100$ molecules and adds the carbon number $n_C$ as a covariate, so
that each descriptor is judged on branching beyond size. Both use the
non-truncating design of Section~\ref{sec:ablation}, with the parent searched
over $(\alpha,\beta) \in (-12,12)^2$ and the Loyola model over the same
box with $\gamma \in (-2,2)$, at $2000$ candidates and seeds
$0$--$49$. The pre-registered criterion for a $\gamma$ effect
required at least $40$ wins in $50$ seeds and a molecule-bootstrap
interval excluding zero (Table~\ref{tab:bp}). The rule is one-sided,
so a negative difference also returns ``no evidence'' and the rule never returns evidence
that $\gamma$ hurts.

\begin{table}[htbp]
\centering
\caption{Measured boiling points (nested leave-one-out $Q^2$). Test A:
decane cohort, refitted within folds. Test B: pooled C6--C10 with the carbon number
$n_C$ in every model (Class., $IRLA$,
LO$_1$: $n_C$ plus that descriptor). Base: training-fold mean (Test A)
or $n_C$ alone (Test B).
Class.: best classical index chosen in each fold. LO$_1$: fixed
$LO(0,0,1)$. Ridge: ridge regression on the degree-pair counts. GM,
LO: medians over seeds $0$--$49$ in the non-truncating boxes.
$\Delta Q^2$: median of LO minus GM. Last column: seeds (of $50$) in
which LO beats GM.}
\label{tab:bp}
\scriptsize
\setlength{\tabcolsep}{1.5pt}
\begin{tabular}{lrrrrrrrrr}
\toprule
Test & Base & Class. & $IRLA$ & LO$_1$ & Ridge & GM & LO & $\Delta Q^2$ & $>$\\
\midrule
decanes ($n = 34$) & $-0.062$ & $0.176$ & $0.167$ & $0.149$ & $0.860$ & $0.514$ & $0.406$ & $-0.072$ & $15$\\
C6--C10, with $n_C$ ($n = 100$) & $0.938$ & $0.967$ & $0.963$ & $0.963$ & $0.991$ & $0.977$ & $0.977$ & $-0.000$ & $25$\\
\bottomrule
\end{tabular}
\end{table}

Both pre-registered verdicts are negative.
On the decane cohort the median $\Delta Q^2 = -0.072$, LO is ahead
in $15$ of $50$ seeds, and the $95\%$ interval is $[-0.156, +0.001]$. In the
pooled size-adjusted set the values are $-0.000$, $25$ of $50$ and $[-0.003, +0.002]$.
In the original box of half-width $2$, reported as the pre-registration
requires, the two tuned models are level on the decanes (medians $0.469$
and $0.465$). The Loyola model is marginally ahead on the pooled set
($0.980$ against $0.976$). Therefore the decane verdict does not depend on the
box. The lower decane $Q^2$ of the hybrid control ($0.406$ against
$0.465$ in the narrow box) reflects the cost of a three-dimensional
search over a wide box. It is not a penalty for $\gamma$.
The pooled cohort was corrected after an external audit found one
nonane's second NIST entry skipped (Deviation D5 in the
repository), and pooling grouped pages also moved one heptane value by
$0.2\,^{\circ}$C. On the uncorrected cohort of $101$ molecules the same test
gave $-0.000$, $24$ of $50$, $[-0.003, +0.002]$, so the correction
changes no conclusion. The intervals are molecule bootstraps
conditional on the fitted predictions, as in Section~\ref{sec:external}.
Ridge regression on the ten degree-pair counts
reaches $Q^2 = 0.860$ on the decanes, where every single index stays
between $0.15$ and $0.51$, and $0.991$ on the pooled set. A permutation
check on the decanes ($200$ random permutations of the boiling points)
gives it a median $Q^2$ of $-0.09$, so the gap is not an artefact of the
cross-validation procedure itself. A permutation test cannot, however,
exclude every form of leakage. Branching effects on the boiling point
are thus spread over several degree-pair counts in a way that no
one-dimensional summary of the family, with or without $\gamma$,
recovers. These conclusions concern acyclic alkanes and are the
primary predictive evidence of this paper, while the
vaporization-enthalpy analyses of Sections~\ref{sec:ablation}
and~\ref{sec:external} are secondary diagnostics. Data, rule,
exclusions and scripts are in the repository.

Data, scripts and verifiers that recompute every tabulated value are
in the reproducibility package~\cite{LoyolaRepo}, where the nonane and
boiling-point protocols were committed before the corresponding
models were fitted. The termwise bounds, the convexity statement, the
independence of the imbalance direction and the order-thirteen profile
collision have also been machine-checked in Lean~4 with Mathlib. There
graphs are modelled by their multisets of edge degree pairs, while equality
characterisations were not formalised. The scope of the formalisation
is described in the Supplementary Material.

% =====================================================================
\section{Conclusion}\label{sec:conclusion}
% =====================================================================

We extended the Gutman--Milovanovi\'c family by an exponential factor
in the normalised degree imbalance of each edge and asked what this
third coordinate adds. On the mathematical side we showed that the imbalance direction is
affinely independent of the two parent directions once a vertex of
degree three is present. The logarithm of the index is globally
convex. The family factorises
into a size coordinate and a shape coordinate. On the resulting ratio
plane the pure-imbalance points discriminate exactly the histogram of
edge imbalances. For trees of small order this histogram carries the
full degree-pair information, so these points attain the discrimination
limit of the whole degree-based class. From order thirteen onward
(sixteen for chemical trees) they provably lose information, and we
exhibited the colliding families. We also bounded the number of order
exchanges between two chemical graphs along the imbalance axis. For
each graph we characterised the parameter directions that leave the
index unchanged up to scale.

On the chemometric side the evidence was negative. An apparent gain
from the imbalance parameter on vaporization enthalpies replicated
under a pre-registered protocol, but it disappeared once the parent
family was searched over a box that did not truncate it.
Pre-registered tests on measured boiling points of decanes and of a
pooled set of C6--C10 alkanes showed no benefit from the third
coordinate. There a linear model on the degree-pair counts
outperformed every single index. The untuned pure-imbalance point is highly
correlated with the parameter-free irregularity index and matched the
tuned models on vaporization enthalpy. The compactness of the family
therefore comes from linear imbalance weighting rather than from
tuning. The third coordinate remains useful as an interpretable direction on
the parent family along which the effect of degree imbalance can be
analysed exactly. Open problems are listed in the Supplementary
Material.

\acknowledgment{The authors thank colleagues for helpful discussions.}

\singlespacing

 \begin{thebibliography}{00}


  \bibitem{AliGutmanFurtula2026BID} A. Ali, I. Gutman, B. Furtula,
  Extremal results on bond incident degree indices of graphs: a survey,
  \textit{Aequationes Math.\/} \textbf{100} (2026) \#66.

  \bibitem{AliGutmanFurtulaAlbalahiHamza2025} A. Ali, I. Gutman,
  B. Furtula, A. M. Albalahi, A. E. Hamza, On chemical and mathematical
  characteristics of generalized degree-based molecular descriptors,
  \textit{AIMS Math.\/} \textbf{10} (2025) 6788--6804.


  \bibitem{Baker1975} A. Baker, \textit{Transcendental Number Theory\/},
  Cambridge Univ. Press, Cambridge, 1975.

  \bibitem{BarmanDas2026HSO} J. Barman, S. Das, Geometric approach to
  degree-based topological index: Hyperbolic Sombor index, \textit{MATCH
  Commun. Math. Comput. Chem.\/} \textbf{95} (2026) 63--94.

  \bibitem{BollobasErdos1998}
  B. Bollob\'as, P. Erd\H{o}s, Graphs of extremal weights, \textit{Ars Comb.\/}
  \textbf{50} (1998) 225--233.

  \bibitem{Brown1986} L. D. Brown, \textit{Fundamentals of Statistical
  Exponential Families with Applications in Statistical Decision
  Theory\/}, IMS Lecture Notes--Monograph Series 9, Institute of
  Mathematical Statistics, Hayward, 1986.


  \bibitem{DasMondal2024ExpGA} K. C. Das, S. Mondal, On exponential
  geometric-arithmetic index of graphs, \textit{J. Math. Chem.\/}
  \textbf{62} (2024) 2740--2760.

  \bibitem{DasMondalRazaPal2025ExpISDD} P. Das, S. Mondal, Z. Raza,
  A. Pal, On exponential inverse symmetric division deg index of graphs,
  \textit{J. Appl. Math. Comput.\/} \textbf{71} (Suppl.~1) (2025) 537--561.


  \bibitem{FurtulaGutmanDehmer2013SS} B. Furtula, I. Gutman, M. Dehmer, On
  structure-sensitivity of degree-based topological indices, \textit{Appl.
  Math. Comput.\/} \textbf{219} (2013) 8973--8978.

  \bibitem{GaoAMC2024} W. Gao, Y. Gao, The extremal trees for
  exponential vertex-degree-based topological indices, \textit{Appl.
  Math. Comput.\/} \textbf{472} (2024) \#128634.

  \bibitem{GranadosJMC2025GM} A. Granados, A. Portilla, Y. Quintana,
  E. Tour\'is, Bounds for the Gutman--Milovanovi\'c index and some
  applications, \textit{J. Math. Chem.\/} \textbf{63} (2025)
  406--418.


  \bibitem{Gutman2026SomborSurvey} I. Gutman, Survey of Sombor-type
  degree-based topological indices, \textit{MATCH Commun. Math. Comput.
  Chem.\/} \textbf{96} (2026) 819--837.

  \bibitem{GutmanBuragohainNyauli2026} I. Gutman, J. Buragohain,
  G. G. Nyauli, Concealed relations between vertex-degree-based
  topological indices, \textit{J. Math. Chem.\/} \textbf{64} (2026) \#39.


  \bibitem{GutmanMilovanovic2020AKCE} I. Gutman, E. Milovanovi\'c,
  I. Milovanovi\'c, Beyond the Zagreb indices, \textit{AKCE Int. J.
  Graphs Comb.\/} \textbf{17} (2020) 74--85.


  \bibitem{GutmanTrinajstic1972} I. Gutman, N. Trinajsti\'c, Graph theory
  and molecular orbitals. Total $\pi$-electron energy of alternant
  hydrocarbons, \textit{Chem. Phys. Lett.\/} \textbf{17} (1972) 535--538.

  \bibitem{HuMATCH2022} Z. Hu, L. Li, X. Li, D. Peng, Extremal graphs
  for topological index defined by a degree-based edge-weight
  function, \textit{MATCH Commun. Math. Comput. Chem.\/} \textbf{88}
  (2022) 505--520.

  \bibitem{Jameson2006} G. J. O. Jameson, Counting zeros of generalised
  polynomials: Descartes' rule of signs and Laguerre's extensions,
  \textit{Math. Gaz.\/} \textbf{90} (2006) 223--234.


  \bibitem{Konstantinova1996Degeneracy} E. V. Konstantinova, The
  discrimination ability of some topological and information distance indices
  for graphs of unbranched hexagonal systems, \textit{J. Chem. Inf. Comput.
  Sci.\/} \textbf{36} (1996) 54--57.


  \bibitem{LabaufGreenshieldsRossini1961} A. Labbauf, J. B. Greenshields,
  F. D. Rossini, Heats of formation, combustion, and vaporization of the
  35 nonanes and 75 decanes, \textit{J. Chem. Eng. Data\/} \textbf{6}
  (1961) 261--263.

  \bibitem{NISTWebBook} P. J. Linstrom, W. G. Mallard (Eds.), \textit{NIST
  Chemistry WebBook, NIST Standard Reference Database Number 69\/}, National
  Institute of Standards and Technology, Gaithersburg, 2025,
  https://doi.org/10.18434/T4D303.


  \bibitem{MolinaSigarreta2025GM} E. D. Molina, J. M.
  Rodr\'iguez-Garc\'ia, J. M. Sigarreta, S. J. Torralbas Fitz, On the
  Gutman--Milovanovi\'c index and chemical applications, \textit{AIMS
  Math.\/} \textbf{10} (2025) 1998--2020.


  \bibitem{MovahediGutmanRedzepovicFurtula2026} F. Movahedi, I. Gutman, I.
  Red\v{z}epovi\'c, B. Furtula, Diminished Sombor index, \textit{MATCH Commun.
  Math. Comput. Chem.\/} \textbf{95} (2026) 141--162.

  \bibitem{LoyolaRepo} A. Natarajan, G. Shiwakoti, M. Arockiaraj,
  The Loyola index family: data, scripts and verifiers,
  https://github.com/Singati2/LOYOLA-PAPER (2026).

  \bibitem{NyauliBuragohain2027DRSO} G. G. Nyauli, J. Buragohain,
  Geometric approach to degree-based topological indices: degree-ratio
  Sombor indices of graphs, \textit{MATCH Commun. Math. Comput. Chem.\/}
  \textbf{97} (2027) 135--164.

  \bibitem{PortillaJMC2025GM} A. Portilla, J. M. Rodr\'iguez,
  I. Salas Lorenzo, J. M. Sigarreta, The Gutman--Milovanovi\'c index
  in mathematical chemistry, \textit{J. Math. Chem.\/} \textbf{63}
  (2025) 2068--2087.

  \bibitem{Rada2019ExpVDB} J. Rada, Exponential vertex-degree-based
  topological indices and discrimination, \textit{MATCH Commun. Math.
  Comput. Chem.\/} \textbf{82} (2019) 29--41.

  \bibitem{Rada2026LinearVDB} J. Rada, A finite-dimensional linear
  framework for vertex-degree-based topological indices, \textit{MATCH
  Commun. Math. Comput. Chem.\/} \textbf{96} (2026) 841--863.

  \bibitem{RakicFurtula2019SS} M. Raki\'c, B. Furtula, A novel method for
  measuring the structure sensitivity of molecular descriptors, \textit{J.
  Chemometrics\/} \textbf{33} (2019) \#e3138.

  \bibitem{Randic1975} M. Randi\'c, Characterization of molecular
  branching, \textit{J. Am. Chem. Soc.\/} \textbf{97} (1975) 6609--6615.

  \bibitem{RetiSharafdini2018IRLA} T. R\'eti, R. Sharafdini,
  \'A. Dr\'egelyi-Kiss, H. Haghbin, Graph irregularity indices used as
  molecular descriptors in QSPR studies, \textit{MATCH Commun. Math.
  Comput. Chem.\/} \textbf{79} (2018) 509--524.
  \bibitem{APIProject44Circ461} F. D. Rossini, K. S. Pitzer, W. J. Taylor,
  J. P. Ebert, J. E. Kilpatrick, C. W. Beckett, M. G. Williams, H. G. Werner,
  \textit{Selected Values of Properties of Hydrocarbons\/}, Natl. Bur.
  Stand. Circ. 461, U.S. Government Printing Office, Washington, 1947.


  \bibitem{ShirdelRezapourSayadi2013HM} G. H. Shirdel, H. Rezapour,
  A. M. Sayadi, The hyper-Zagreb index of graph operations,
  \textit{Iran. J. Math. Chem.\/} \textbf{4} (2013) 213--220.


\bibitem{Sigarreta2022ExpVDB} J. M. Sigarreta, Extremal problems on
  exponential vertex-degree-based topological indices, \textit{Math.
  Biosci. Eng.\/} \textbf{19} (2022) 6985--6995.

  \bibitem{VukicevicFurtula2009GA} D. Vuki\v{c}evi\'c, B. Furtula, Topological
  index based on the ratios of geometrical and arithmetical means of
  end-vertex degrees of edges, \textit{J. Math. Chem.\/} \textbf{46} (2009)
  1369--1376.


  \bibitem{WilhoitZwolinski1971} R. C. Wilhoit, B. J. Zwolinski,
  \textit{Handbook of Vapor Pressures and Heats of Vaporization of
  Hydrocarbons and Related Compounds\/}, Texas A\&M Research Foundation,
  College Station, 1971.

  \bibitem{ZhouTrinajstic2009sumconn}
  B. Zhou, N. Trinajsti\'c, On a novel connectivity index, \textit{J. Math.
  Chem.\/} \textbf{46} (2009) 1252--1270.

  \bibitem{ZhouTrinajstic2010} B. Zhou, N. Trinajsti\'c, On general
  sum-connectivity index, \textit{J. Math. Chem.\/} \textbf{47} (2010)
  210--218.


\end{thebibliography}

\end{document}

